\chapter{Boolean Analysis — InputMoments}
%%% generated by the blueprint pipeline from TCSlib.BooleanAnalysis.GaussianSpace.InputMoments
\section{Overview}
Standard Gaussian variables and centered random signs satisfy the moment conditions used in fourth-order invariance.

\section{Declarations}

\begin{lemma}[Fourth moment of a standard Gaussian]
\lean{BooleanAnalysis.GaussianSpace.integral_fourth_standardGaussian}
\leanok
\uses{BooleanAnalysis.ThresholdFunctions.standardGaussianMeasure}
\difficulty{6}
If $Z$ is a standard real Gaussian random variable, then $\mathbb E[Z^4]=3$.
\end{lemma}

\begin{lemma}[Vanishing third Gaussian moment]
\lean{BooleanAnalysis.GaussianSpace.integral_third_standardGaussian}
\leanok
\uses{BooleanAnalysis.ThresholdFunctions.standardGaussianMeasure}
\difficulty{4}
If $Z$ is a standard real Gaussian random variable, then $\mathbb E[Z^3]=0$.
\end{lemma}

\begin{lemma}[Second moment of a standard Gaussian]
\lean{BooleanAnalysis.GaussianSpace.integral_second_standardGaussian}
\leanok
\uses{BooleanAnalysis.ThresholdFunctions.standardGaussianMeasure}
\difficulty{3}
If $Z$ is a standard real Gaussian random variable, then $\mathbb E[Z^2]=1$.
\end{lemma}

\begin{lemma}[Gaussian invariance input moments]
\lean{BooleanAnalysis.GaussianSpace.hasInvarianceMoments_standardGaussian}
\leanok
\uses{BooleanAnalysis.GaussianSpace.HasInvarianceMoments, BooleanAnalysis.GaussianSpace.integral_fourth_standardGaussian, BooleanAnalysis.GaussianSpace.integral_second_standardGaussian, BooleanAnalysis.GaussianSpace.integral_third_standardGaussian, BooleanAnalysis.ThresholdFunctions.standardGaussianMeasure}
\difficulty{4}
A standard real Gaussian random variable satisfies the invariance moment conditions: it has finite fourth absolute moment, its first three moments are $0,1,0$, and its fourth moment is at most $9$.
\end{lemma}

\begin{lemma}[Centered sign input moments]
\lean{BooleanAnalysis.GaussianSpace.hasInvarianceMoments_of_sq_eq_one}
\leanok
\uses{BooleanAnalysis.GaussianSpace.HasInvarianceMoments}
\difficulty{4}
Let $X$ be a real random variable on a probability space, measurable up to a null set, with $\mathbb EX=0$ and $X(\omega)^2=1$ for every $\omega$. Then $X$ has finite fourth absolute moment and satisfies $\mathbb EX^2=1$, $\mathbb EX^3=0$, and $\mathbb EX^4\le9$, so it satisfies the invariance moment conditions.
\end{lemma}

\begin{lemma}[Uniform sign input moments]
\lean{BooleanAnalysis.GaussianSpace.hasInvarianceMoments_uniformBool}
\leanok
\uses{BooleanAnalysis.GaussianSpace.HasInvarianceMoments, BooleanAnalysis.GaussianSpace.hasInvarianceMoments_of_sq_eq_one, BooleanAnalysis.boolToSign}
\difficulty{3}
Let $B$ be a uniform bit and put $X=1$ when $B$ is false and $X=-1$ when $B$ is true. Then $X$ satisfies the invariance moment conditions, including finite fourth absolute moment, first three moments $0,1,0$, and fourth moment at most $9$.
\end{lemma}
